\apendice{Descripción de adquisición y tratamiento de datos} 
En este aparatado se detallará la obtención de datos, desde los solicitados mediante APIs hasta los de elaboración propia.


\section{Descripción formal de los datos}
El proceso de obtención se puede dividir en tres bloques:
\begin{itemize}
    \item \textbf{Los datos obtenidos mediante la API de la AEMET}: para este proceso se desarrollo una función llamada \textit{actualizar\_datos\_climaticos()} que encapsula todo el proceso de descarga, transformación y almacenamiento de datos meteorológicos en un bloque único.
    El proceso comienza con la consrucción del \textit{endpoint} disponible en la \href{https://opendata.aemet.es/dist/index.html?}{AEMET} y la petición de datos. 
    \begin{figure}[H]
        \centering
        \includegraphics[width=0.8\linewidth]{img/peticion.png}
        \caption{Petición}
        \label{fig:peticion}
    \end{figure}
    Para temperaturas, se ha escogido el \textit{endpoint} que proporciona observaciones convencionales como humedad, precipitación, viento y entre otras se encuentra la temperatura de cada estación meteorológica que se registra cada hora. La respuesta no son los datos directamente, sino que se trata de un \textit{JSON} con información sobre el estado de la solicitud, información sobre los datos proporcionados y el enlace a los datos obtenidos, figura \ref{fig:RESPUESTA}.
    \begin{figure}[H]
        \centering
        \includegraphics[width=0.8\linewidth]{respuesta.png}
        \caption{Estructura de la respuesta de la API.}
        \label{fig:RESPUESTA}
    \end{figure}
    Se procede a extraer los datos realizando una descarga del contenido del enlace proporcionado, conservando su estructura original, posteriormente realizando una conversión total al estándar de codificación \textit{UTF-8} para mejor gestión de datos y evitar errores.
    \begin{figure}[H]
        \centering
        \includegraphics[width=0.8\linewidth]{img/json.png}
        \caption{Datos observacionales. }
        \label{fig:observaciones}
    \end{figure}
    Como se puede observar en la figura \ref{fig:observaciones}, hay varios registros de distintas variables pero solo nos interesan los atributos \textit{idema}, \textit{tamax} y \textit{tamin} para el proyecto. Los registros son por horas y por \textit{idemas}, que son los identificadores de las estaciones meteorológicas que corresponden a una provincia, por ello se realiza un filtro de la temperatura mínima y máxima de todo el día de cada identificador.
    Tener los datos identificados por idemas suponía un problema a la hora de representarlos geográficamente, por ello se acudió a la funcion \textit{aemet\_stations()} para obtener las correspondencias de cada \textit{idema} a su respectiva provincia, para realizar el proceso de asociación \textit{idema}-->provincia. Finalmente se agrupan todas los idemas en su provincia con su temperatura promedio máxima y mínima junto a la fecha de adquisición.
    \begin{figure}[H]
        \centering
        \includegraphics[width=0.6\linewidth]{img/ideme.png}
        \caption{Provincias de los idemas.}
        \label{fig:idemas}
    \end{figure}
    De manera análoga, se obtuvo los datos de predicción de índice ultravioleta, cambiando el \textit{endpoint} por el de predicción índice ultravioleta con código 0 para tener los datos del día actual. A diferencia de las temperaturas, estos datos no se necesitaba su asociación \textit{idema}-->provincia, ya que estos datos venían dados en provincias.
    \begin{figure}[H]
        \centering
        \includegraphics[width=0.8\linewidth]{img/peticiouv.png}
        \caption{Petición índice UV}
        \label{fig:uvpet}
    \end{figure}
    \begin{figure}[H]
        \centering
        \includegraphics[width=0.8\linewidth]{img/respu.png}
        \caption{Datos índice UV.}
        \label{fig:enter-label}
    \end{figure}
    Para la obtención de altitudes, se extrajeron de la función  \textit{aemet\_stations()} agrupando los idemas en su provincia correspondiente calculando el promedio de altitud del conjunto de idemas de cada provincia.
    Como paso final se implementó un diccionario de equivalencias de los nombres de las provincias, ya que cada variable presentaba un nombre diferente ya sea por mayúsculas, acentos o segundos idiomas de la provincia, para un mejor manejo a la hora de representarlos en el mapa.
    
    \item \textbf{Los datos de las defunciones por melanoma maligno de la piel}: se obtuvieron del paquete de la INE, \textit{INEbaseR}, que conecta con la API del Instituto Nacional de Estadística. Se filtran solo los de melanoma maligno de la piel, y de nuevo se realiza la transformación a un nombre común de las provincias para que corresponda con los valores de las provincias de la representación geográfica.
    
    \item \textbf{Los datos de elaboración propia:} estos datos han sido los de la densidad de población de cada provincia para calcular la tasa de incidencia del melanoma \cite{AnexoProvinciasCiudades2025} .
    
\end{itemize}
